% answers_short.tex —— 2000 年·上海卷·文科 · 简短答案速查
% 由 tools/make_short_answers.py 自 answers.tex 的 \gksolhead/\gkanswer 提取，
% 并与 2000/2000_数学真题_上海文_答案要点.md 逐题比对后生成，勿手改（改源 MD 后重跑）。

\begin{gknotice}
\textbf{使用说明}：本节只给最终答案便于快速核对，完整推导见解析版。答案为编写者独立推导，\textbf{非官方标准答案}。
\end{gknotice}

\gkpart{填空题}
\begin{itemize}[leftmargin=2.6em]
  \item[\textbf{1.}] $m=\dfrac{3}{2}$
  \item[\textbf{2.}] $\left(\dfrac{1}{2},3\right)$
  \item[\textbf{3.}] $(-4,0)$、$(6,0)$
  \item[\textbf{4.}] $e^{-2}$
  \item[\textbf{5.}] $b=1$
  \item[\textbf{6.}] $9$ 年
  \item[\textbf{7.}] 例如三条侧棱长相等
  \item[\textbf{8.}] $f(x)=x$（$1\leq x\leq2$）
  \item[\textbf{9.}] $-462$
  \item[\textbf{10.}] $\dfrac{1}{14}$
  \item[\textbf{11.}] $(0,5)$
  \item[\textbf{12.}] $b_1b_2\cdots b_n=b_1b_2\cdots b_{17-n}$（$1\leq n<17$）
\end{itemize}

\gkpart{选择题}
\begin{center}\begin{tabular}{ccccc}
\toprule
题号 & 13 & 14 & 15 & 16 \\
\midrule
答案 & C & A & A & D \\
\bottomrule
\end{tabular}\end{center}

\gkpart{解答题}
\begin{itemize}[leftmargin=2.6em]
  \item[\textbf{17.}] $\left(-\dfrac{9}{5},\dfrac{1}{5}\right)$
  \item[\textbf{18.}] $\dfrac{8}{3}$
  \item[\textbf{19.}] (1) $\dfrac{7}{2}$（在 $x=1$ 处取得）；(2) $a>-3$
  \item[\textbf{20.}] (1) $(4\sqrt{2},45^\circ)$；(2) 在 $(7,0)$ 截住；截住指令为 $(5.00,-98.13^\circ)$
  \item[\textbf{21.}] (1) U/PENDING；按固定印式字面读得 $b_n=2000\left(\dfrac{a}{10}\right)^2$；(2) U/PENDING；按固定印式字面读法，$0<a<10$ 均满足；(3) U/PENDING；按固定印式字面读法不存在最大有限项数
  \item[\textbf{22.}] (1) $m=\sqrt{3}$，$x'=x+\sqrt{3}y$，$y'=\sqrt{3}x-y$；(2) $\left(\dfrac{3\sqrt{3}}{4},\dfrac{1}{4}\right)$；(3) $k=\dfrac{1}{\sqrt{3}}$ 或 $k=-\sqrt{3}$
\end{itemize}
